Las conclusiones constituyen otro capítulo. Lo principal es mencionar si se aceptó o rechazó la hipótesis y si se cumplieron los objetivos. Además, se indica cómo contribuye el trabajo a resolver el problema original. Pueden escribirse en forma de párrafo o con viñetas. Para redactarlas, tome en cuenta lo siguiente:  

\begin{itemize}
    \item Son específicas, no generalizadas. 

    \item Se escriben en orden.

    \item Evitan expresiones como «yo pienso», «opino», «desde mi experiencia», etc.; el investigador debe desaparecer, por lo que se redacta en tercera persona.
    
\end{itemize}